\section{Extended Formalization}
\label{app:formalization}

\paragraph{Harness and residual knowledge.}
The fixed harness $H$ contains the executor implementation, system prompt,
available tools, environment interface, and retrieval mechanism. In our
implementation, $H$ is instantiated by \texttt{pi-agent}, and its action space
contains only bash commands. We use \emph{residual procedural knowledge} for
any reusable condition, action, implementation detail, or error path that
improves execution but is absent from $H$. This definition is operational: a
knowledge item is residual if it can be learned from a rollout and changes
later task execution without changing the model or harness.

\paragraph{Round update.}
For the $n=127$ benchmark tasks $Q=\{q_i\}_{i=1}^{n}$, the execute phase under $M_r$
produces rollout batch $B_r$, the set $\{\tau_i^r\}_{i=1}^{n}$ in which
$\tau_i^r$ is the trace of the attempt at $q_i$ during round $r$. A feedback function $E$ maps each rollout and its
resulting environment state to $f_i^r$. Scope $S$ selects passed checkpoint
evidence only or both passed and failed checkpoint evidence. The consolidate
phase applies
\[
M_{r+1}=C\!\left(M_r,
  \{(\tau_i^r,f_i^r)\mid q_i\in Q\};S\right).
\]
The memory store remains frozen within an execute phase, so all rollouts in one
batch use the same $M_r$; updates become visible only in the next round.

\paragraph{Procedural sufficiency.}
A consolidated procedure should specify (i) its purpose, (ii) conditions under
which it applies, (iii) executable actions and implementation details, and
(iv) recovery behavior when an action fails. Consolidation is successful when
it transforms multiple rollout-specific observations into a procedure that
improves execution on other tasks sharing those conditions. This criterion
permits different surface representations while excluding pure replay as the
target behavior.

\begin{table}[th]
\caption{\textbf{Memory types and evaluation requirements.} Procedural memory
differs because repeated trials should refine a reusable action policy, and its
quality is established by execution rather than recall.}
\label{tab:memory-types}
\centering
\small
\begin{tabularx}{\linewidth}{@{}lXXX@{}}
\toprule
 & Episodic & Semantic & Procedural \\
\midrule
Stored content &
Specific event and context &
General fact or concept &
Condition-action procedure \\
Repeated trials &
Remain distinct events &
Duplicate facts merge &
Refine one reusable procedure \\
Failure value &
Another event &
Wrong facts are discarded &
Reveals missing conditions or error paths \\
Generalization target &
Event recall &
Concept-level recall &
Execution across task instantiations \\
Primary evaluation &
Recall accuracy &
Recall accuracy &
Execution success \\
\bottomrule
\end{tabularx}
\end{table}
